\chapter{Lời hứa có biên được chứng nhận}

Chương này trình bày cơ chế trung tâm của khóa luận. Cấu trúc chương đi từ mô hình bài toán mở
rộng, qua phép phân rã ba mốc và sổ cái cam kết, tới phần lý thuyết về hai cơ sở tính sổ cái và
khung so sánh bốn họ ràng buộc lời hứa, rồi kết thúc bằng kiến trúc cài đặt và cơ chế bảo đảm
tính tái lập.

\section{Tổng quan cơ chế}

Ý tưởng có thể phát biểu trong một đoạn. Hệ thống điều phối, ngoài việc gán khách cho xe như
thông thường, còn \emph{công bố} cho mỗi hành khách một giờ đón và một giờ trả. Mỗi khi hệ
thống quyết định lại và công bố một giờ mới, độ dời giữa giờ mới và giờ cũ được ghi vào một sổ
cái riêng của hành khách đó. Sổ cái cộng dồn theo thời gian và chỉ tăng. Trước khi chấp nhận
một phương án, hệ thống kiểm tra xem phương án ấy có làm sổ cái của bất kỳ hành khách nào vượt
trần $\beta$ hay không; nếu có, phương án bị loại khỏi tập lựa chọn ngay tại đó.

Điểm tinh tế nằm ở chỗ \emph{độ dời nào bị tính vào sổ cái}. Nếu tính toàn bộ độ dời mà hành
khách nhìn thấy, hệ thống sẽ bị tính phí cho cả những thay đổi do xe chạy chậm hoặc do đường
tắc --- những thứ nó không gây ra. Để tránh điều đó, hệ thống dựng thêm một mốc tham chiếu phản
thực: \emph{chính kế hoạch cũ đó, nhưng chiếu lại từ vị trí thật hiện tại của xe}. Hiệu giữa
phương án mới và mốc phản thực này là phần do quyết định gây ra; hiệu giữa mốc phản thực và giờ
đã công bố là phần thay đổi khi giữ nguyên kế hoạch. Chỉ phần thứ nhất bị tính.

\begin{figure}[htbp]
\centering
\begin{tikzpicture}[
  node distance=8mm,
  hop/.style={draw,rounded corners=3pt,align=center,inner sep=5pt,
              font=\small,minimum height=10mm,fill=xamnhat,draw=xamvien},
  mui/.style={-{Latex[length=2.2mm]},thick}
]
\node[hop] (a) {Sự kiện mới\\(yêu cầu, vị trí xe)};
\node[hop,right=of a] (b) {Sinh phương án\\cho từng xe};
\node[hop,right=of b] (c) {Dựng ba mốc\\$v$, $e$, $p$};
\node[hop,right=of c] (d) {Đối chiếu\\sổ cái với $\beta$};
\node[hop,below=10mm of d] (e) {Chọn phương án\\bằng ILP};
\node[hop,left=of e] (f) {Cập nhật sổ cái\\và công bố lại};
\node[hop,left=of f] (g) {Ghi chứng cứ\\niêm phong};
\draw[mui] (a)--(b); \draw[mui] (b)--(c); \draw[mui] (c)--(d);
\draw[mui] (d)--(e); \draw[mui] (e)--(f); \draw[mui] (f)--(g);
\end{tikzpicture}
\caption{Vòng quyết định của hệ thống. Ba bước ở giữa là phần bổ sung so với kiến trúc điều
phối thông thường.}
\label{hinh:vongquyetdinh}
\end{figure}

\section{Mô hình bài toán mở rộng}

\subsection{Lời hứa công bố}

Mở rộng mô hình ở mục 2.1, ta bổ sung khái niệm lời hứa. Với mỗi hành khách $r$ đã được nhận,
hệ thống duy trì một \emph{lời hứa đang mở}, gồm định danh xe được phân công, điểm đón, điểm
trả, giờ đón dự kiến và giờ trả dự kiến. Lời hứa này được công bố cho hành khách và được ghi
vào biên bản. Mỗi khi hệ thống ra quyết định, lời hứa có thể được công bố lại với giá trị mới.

Trạng thái cam kết của một hành khách đi qua ba giai đoạn: \emph{đã nhận} (đã được xác nhận
nhưng xe chưa tới), \emph{đang chờ đón}, và \emph{đang trên xe}. Việc phân biệt ba giai đoạn
là cần thiết vì các chiều cam kết có ý nghĩa khác nhau ở từng giai đoạn --- chẳng hạn giờ đón
không còn ý nghĩa sau khi khách đã lên xe.

\subsection{Các chiều cam kết}

Hệ thống theo dõi mười chiều cam kết, chia thành ba nhóm. Nhóm thời gian gồm tổng độ dời giờ
đón và tổng độ dời giờ trả, tính bằng mili giây. Nhóm định danh gồm số lần đổi xe, số lần đổi
điểm đón, số lần đổi điểm trả. Nhóm cấu trúc gồm khoảng cách dời điểm đón và điểm trả tính bằng
mili-mét, số lần đảo thứ tự hành khách cũ, số điểm dừng chèn thêm trước khi đón, và số lần sửa
đổi mang tính vật chất.

Mỗi chiều có thể được đặt một trần cứng, hoặc để không giới hạn. Trong thực nghiệm, cấu hình
được dùng nhiều nhất đặt trần cho chiều \emph{tổng độ dời giờ trả} và khóa cứng ba chiều định
danh về không, tức cấm hoàn toàn việc đổi xe hay đổi điểm dừng sau khi đã hứa.

\section{Phép phân rã ba mốc}
\label{muc:bamoc}

\subsection{Định nghĩa}
\label{muc:dinhnghia-bamoc}

Tại một thời điểm quyết định, với hành khách $r$ đang có lời hứa mở trên xe $v$, hệ thống dựng
ba mốc:

\begin{itemize}[leftmargin=1.2cm,itemsep=2pt,topsep=3pt]
\item $v$ --- \textbf{previous}: lời hứa đã công bố lần gần nhất cho hành khách;
\item $e$ --- \textbf{exogenous}: chính kế hoạch đó, chiếu lại từ vị trí thật hiện tại của xe.
Cùng lộ trình, cùng phân công, chỉ khác điểm xuất phát và thời điểm chiếu;
\item $p$ --- \textbf{proposed}: lời hứa mà phương án đang xét sẽ tạo ra.
\end{itemize}

Từ ba mốc này định nghĩa ba đại lượng:

\begin{align}
\text{độ dời do quyết định} \quad & A = |p - e|\\
\text{thay đổi khi giữ tuyến} \quad & d = |e - v|\\
\text{độ dời khách nhìn thấy} \quad & B = |p - v|
\end{align}

Mốc $e$ là mấu chốt. Nó trả lời câu hỏi phản thực ``nếu nhà điều hành \emph{không làm gì}, giờ
hẹn của hành khách này bây giờ sẽ là mấy?''. Chênh lệch giữa câu trả lời đó và giờ đã hứa là
phần thay đổi khi giữ nguyên kế hoạch --- gồm tắc đường nhưng cũng gồm tiến độ thực và độ trễ thực
thi của xe (xem đoạn cuối mục \ref{muc:phanra}); chênh lệch giữa phương án đang xét và câu trả
lời đó chính là phần nhà điều hành gây ra.

\paragraph{$A$ là đại lượng đối xứng.} Cài đặt tính $A = |p - e|$ theo đúng trị tuyệt đối
(\texttt{PromiseDeltaCalculator}, hàm \texttt{AbsoluteDifference}): một quyết định kéo lời hứa
\emph{sớm} hơn mốc không-làm-gì cũng bị tính phí bằng đúng độ lệch đó, giống hệt một quyết định
đẩy muộn. Không có quy ước một chiều nào trong mã.

Hệ quả cần nói thẳng là cơ sở hành vi \emph{có thể} phạt một sửa chữa có lợi cho khách. Xét
$v = 600$, $e = 720$, $p = 600$ (giây): giao thông đẩy lời hứa muộn đi 120 giây, và phương án
đang xét kéo nó về đúng giờ đã hứa. Khi đó $A = 120$ trong khi $B = 0$ --- khách không thấy gì
thay đổi, nhưng sổ cái hành vi vẫn bị trừ 120 giây. Một trần hành vi chặt có thể loại đúng phương
án mà khách muốn nhất.

Trên panel thực nghiệm, trường hợp này \emph{chưa từng xảy ra}: trong cả $479$ lần công bố có
$A > 0$ ở nhánh B1 và $11$ lần ở nhánh C1, dấu của $p - e$ xác định được hoàn toàn và
\emph{luôn dương}. Nguyên nhân mang tính cơ học --- bộ sinh phương án chỉ chèn thêm điểm dừng,
nên chỉ có thể đẩy lùi giờ của khách cũ chứ không kéo sớm lên được. Vì vậy trên hệ hiện tại
$A = p - e$ trong mọi quan sát. Nhưng đó là tính chất \emph{của bộ sinh}, không phải của cơ sở
hành vi; với một bộ sinh biết sắp xếp lại lộ trình, giới hạn nêu trên sẽ trở nên thực tế.

\subsection{Ví dụ minh họa}

Một hành khách đặt xe lúc 9:00 và được báo sẽ được đón lúc 9:10. Đến 9:05, hệ thống cân nhắc
thay đổi kế hoạch vì có khách mới. Giả sử tại thời điểm đó xe đang chậm hơn dự kiến 40 giây.

\begin{table}[htbp]
\centering
\caption{Ba mốc trong một tình huống cụ thể}
\label{bang:bamoc}
\small
\begin{tabular}{@{}llr@{}}
\toprule
\textbf{Mốc} & \textbf{Ý nghĩa} & \textbf{Giá trị}\\
\midrule
$v$ & Giờ đã hứa với khách & 9:10:00\\
$e$ & Giờ sẽ đến nếu không làm gì (xe đang chậm 40 s) & 9:10:40\\
$p$ & Giờ sẽ đến nếu đổi kế hoạch & 9:11:10\\
\bottomrule
\end{tabular}
\end{table}

Với phương án đổi kế hoạch, cơ sở hành vi tính $9{:}11{:}10 - 9{:}10{:}40 = 30$ giây, còn cơ sở
kết quả tính $9{:}11{:}10 - 9{:}10{:}00 = 70$ giây. Khác biệt 40 giây chính là phần xe chạy chậm.

Bây giờ xét phương án \emph{không làm gì}, tức $p = e = 9{:}10{:}40$. Cơ sở hành vi tính 0 giây,
nên phương án này luôn hợp lệ. Cơ sở kết quả tính 40 giây --- đã vượt trần 30 giây. Theo cách
tính thứ hai, \emph{ngay cả việc không làm gì cũng là vi phạm}. Nếu mọi phương án khác cũng vượt,
chiếc xe không còn kế hoạch hợp lệ nào và hệ thống dừng. Đây là hạt nhân của toàn bộ phần lý
thuyết ở mục \ref{muc:haicosotinh}.

\section{Sổ cái cam kết}

Với mỗi cặp (hành khách, chiều cam kết), hệ thống duy trì một bộ đếm cộng dồn. Mỗi lần công bố
lại lời hứa, bộ đếm được cập nhật

\begin{equation}
L_{r,\text{sau}} = L_{r,\text{trước}} + \delta_r
\end{equation}

trong đó $\delta_r$ là độ dời của lần công bố ấy theo cơ sở đang dùng. Vì $\delta_r \ge 0$, bộ
đếm đơn điệu tăng.

Một phương án được coi là hợp lệ về mặt cam kết nếu với mọi hành khách $r$ bị ảnh hưởng và mọi
chiều $k$ có đặt trần, ta có $L_{r,k,\text{sau}} \le \beta_k$. Phương án vi phạm bị loại khỏi
tập lựa chọn trước khi bài toán quy hoạch nguyên được giải, chứ không bị phạt trong hàm mục
tiêu.

Ba đặc điểm của sổ cái này cần được nhấn mạnh vì chúng là chỗ khác biệt so với văn liệu. Thứ
nhất, nó \emph{theo từng hành khách}: không có phép cộng gộp nào làm loãng người chịu thiệt
nặng nhất. Thứ hai, nó \emph{cộng dồn theo thời gian}: một hành khách bị dời năm lần, mỗi lần
mười giây, tiêu tốn đúng bằng một hành khách bị dời một lần năm mươi giây. Thứ ba, nó là
\emph{ràng buộc cứng} chứ không phải số hạng mục tiêu, nên lời hứa cưỡng chế được chứ không
chỉ là chỉ tiêu.

\section{Hai cơ sở tính, các định lý và khung bốn họ cổng}
\label{muc:haicosotinh}

\subsection{Phát biểu hai cơ sở}

Từ mục \ref{muc:bamoc}, có đúng hai cách tự nhiên để tính sổ cái, và chúng không tương đương.

\begin{dinhnghia}[Cơ sở hành vi]
Sổ cái cộng dồn đại lượng $A = |p-e|$: ``tôi làm bạn trễ thêm bao nhiêu so với lẽ ra bạn đã bị''.
Trong cài đặt, đây là \texttt{CommitmentBudgetBasis.DecisionInduced}. Tương ứng với nghĩa vụ
phương tiện trong mục 2.4.
\end{dinhnghia}

\begin{dinhnghia}[Cơ sở kết quả]
Sổ cái cộng dồn đại lượng $B = |p-v|$: ``bạn trễ bao nhiêu so với điều tôi đã nói''. Trong cài
đặt, đây là \texttt{CommitmentBudgetBasis.CustomerVisible}. Tương ứng với nghĩa vụ kết quả.
\end{dinhnghia}

\subsection{Bất biến nền: phương án không làm gì luôn tồn tại}

Toàn bộ phần lý thuyết dựa trên một bất biến của bộ sinh phương án: tập phương án của mỗi xe
\emph{luôn chứa} phương án giữ nguyên kế hoạch hiện tại. Bất biến này được khẳng định độc lập ở
ba vị trí trong mã nguồn --- bộ chọn đội xe ứng viên, bộ giữ lại danh mục phương án, và bộ đánh
giá phương án --- và nó cũng bắt buộc về mặt mô hình: khác với ràng buộc (2.3) dạng ``không quá
một'' của kiến trúc gốc, bài toán chọn của hệ thống này buộc mỗi xe chọn \emph{đúng} một phương
án, trong đó có đúng một phương án không làm gì.

Với phương án không làm gì, kế hoạch không đổi, nên phép chiếu của nó từ vị trí hiện tại
\emph{chính là} mốc $e$. Tức là $p(\text{no-op}) = e$.

\subsection{Phân rã cộng tính và độ khít của nó}
\label{muc:phanra}

Ba đại lượng liên hệ với nhau qua bất đẳng thức tam giác:

\begin{equation}
\boxed{\,B \le A + d\,}
\label{eq:phanra}
\end{equation}

Dấu bằng xảy ra khi và chỉ khi hai hiệu $(p-e)$ và $(e-v)$ không trái dấu, tức khi ba mốc xếp
cùng một chiều. Khóa luận này \emph{không} khẳng định điều kiện đó luôn đúng. Nó đo xem điều kiện
đó đúng đến mức nào, và chỉ ra rằng phần sai lệch nhỏ hơn mọi ngưỡng diễn giải được dùng về sau.

\paragraph{Vì sao dấu bằng có thể hỏng.} Cả ba đại lượng đều là trị tuyệt đối. Nếu phần thay đổi
khi giữ tuyến kéo $e$ \emph{sớm} hơn $v$ trong khi quyết định đẩy $p$ \emph{muộn} hơn $e$, hai hiệu
ngược dấu triệt tiêu một phần và $|p-v| < |p-e| + |e-v|$. Đây không phải trường hợp hiếm: phần
thay đổi khi giữ tuyến trên panel này gần như đối xứng --- đo được 214.924 lần sớm hơn so với 206.224 lần muộn
hơn trên 421.752 lần sửa --- trong khi phần do quyết định, trên thực nghiệm, luôn cùng một chiều
(luôn làm muộn đi, mục \ref{muc:dinhnghia-bamoc}). Một đại lượng luôn cùng chiều gặp một đại lượng hai chiều thì
triệt tiêu từng phần là điều phải xảy ra.

\paragraph{Độ khít đo được ở mức từng lần công bố.} Bảng \ref{bang:khit-tung-lan} kiểm tra
\eqref{eq:phanra} trên từng bản ghi \texttt{promisePublished} của chiều
\texttt{drop\_eta\_total\_ms}, thay vì chỉ trên tổng.

\begin{table}[htbp]
\centering
\caption{Độ khít của phân rã cộng tính, tính trên từng lần công bố (panel 20 ô)}
\label{bang:khit-tung-lan}
\small
\begin{tabular}{@{}lrr@{}}
\toprule
& \textbf{B1} & \textbf{C1}\\
\midrule
Số lần công bố & 54.311 & 45.542\\
$A + d = B$ đúng chính xác & 54.083 \;($99{,}58\%$) & 45.539 \;($99{,}99\%$)\\
$A + d > B$ (bất đẳng thức chặt) & 228 \;($0{,}42\%$) & 3\\
\midrule
Số lần \emph{cả hai} $A>0$ và $d>0$ & 476 \;($0{,}88\%$) & 11\\
\quad trong đó dấu bằng hỏng & \textbf{228} \;(\textbf{$47{,}9\%$}) & 3 \;($27{,}3\%$)\\
\midrule
Độ hụt khi hỏng: trung vị & 168 ms & ---\\
Độ hụt khi hỏng: cực đại & 1.380 ms & 208 ms\\
\textbf{Độ hụt trên tổng $A+d$} & \textbf{$0{,}0748\%$} & \textbf{$0{,}0059\%$}\\
\bottomrule
\end{tabular}
\end{table}

\paragraph{Cách đọc bảng này cho đúng.} Con số $99{,}58\%$ không phải một phép kiểm định mạnh.
Trong $99{,}12\%$ số lần công bố, \emph{ít nhất một} trong hai đại lượng $A$, $d$ bằng không;
khi $e = v$ thì $B = A$ hiển nhiên, khi $p = e$ thì $B = d$ hiển nhiên. Phép kiểm định thật chỉ
tồn tại ở 476 lần công bố mà cả hai đều khác không --- và ở đó dấu bằng \emph{hỏng gần một nửa
số lần}. Vì vậy phát biểu đúng là bất đẳng thức \eqref{eq:phanra} kèm độ khít đã đo, chứ không
phải một đẳng thức.

\paragraph{Vì sao kết luận quy trách vẫn vững.} Phần hụt cộng dồn là 53.530 ms trên tổng
$71{,}58$ triệu ms, tức $0{,}0748\%$; ở nhánh C1 là $0{,}0059\%$. Sàn nhiễu quỹ đạo của toàn bộ
thực nghiệm là $\pm 0{,}3$ điểm phần trăm. Phần hụt nhỏ hơn ngưỡng đó khoảng hai bậc độ lớn. Do
đó mọi kết luận quy trách ở mục \ref{muc:quytrach} --- kể cả con số $89{,}06\%$ --- không thay
đổi dù ta gán toàn bộ phần hụt về phía bất lợi nhất.

\paragraph{Bằng chứng chính cho phép phân rã không phải là dấu bằng.} Trước hết cần nói rõ $d$
đo cái gì. Mốc $e$ được chiếu trên trạng thái \emph{sau khi đã áp dụng toàn bộ lô sự kiện} của
kỷ nguyên đó (\texttt{CommitmentDecisionValidator} đọc tuyến cũ từ \texttt{ReducedState}), nên
$d$ gồm mọi thứ đã đổi trong lô khi giữ nguyên tuyến: tiến độ thực của xe, độ trễ thực thi, cập
nhật thời gian di chuyển, và các sự kiện vòng đời khác. Vì vậy \emph{không} được gọi $d$ là
``tác động thuần của giao thông''; tên chính xác là \emph{phần thay đổi khi chiếu lại tuyến giữ
nguyên trên trạng thái mới}. Và vì tiến độ thực của xe phụ thuộc tuyến đã chọn, $d$ về nguyên
tắc \emph{có thể} phụ thuộc yếu vào chính sách.

Do đó tính bất biến của $d$ là một \emph{tính chất phải kiểm bằng thực nghiệm}, không phải một
định lý --- và nó là phép kiểm mạnh nhất cho phép phân rã. Nếu phép tách hai thành phần bị sai,
$d$ sẽ trôi theo chính sách. Trên hai nhánh B1 và C1, đại lượng $A$ thay đổi \textbf{507 lần}
(từ 64.863.133 ms xuống 128.020 ms) trong khi $d$ chỉ lệch \textbf{$0{,}37\%$} (6.719.286 ms so
với 6.694.740 ms). Thành phần lẽ ra phải phản ứng với chính sách thì phản ứng; thành phần lẽ ra
gần như bất biến thì gần như bất biến.

\subsection{Ba định lý}
\label{muc:badinhly}

\begin{dinhly}[Cơ sở hành vi không bao giờ loại phương án không làm gì]
\label{dl:hanhvi}
Dưới cơ sở hành vi, ràng buộc cam kết không bao giờ loại phương án không làm gì, nên nó không bao
giờ tự làm rỗng tập phương án của một xe.
\end{dinhly}

\begin{proof}
Với phương án không làm gì, $A = |p - e| = |e - e| = 0$. Do đó sổ cái sau bằng sổ cái trước,
mà sổ cái trước không vượt $\beta$ vì mọi lần công bố trước đó đều đã qua kiểm tra. Vậy ràng buộc
cam kết không loại phương án không làm gì; nếu phương án đó cũng qua mọi kiểm tra không thuộc cam
kết (giả thiết (N) ở mục \ref{muc:hocong}), tập phương án có ít nhất một phần tử.
\end{proof}

Ý nghĩa của định lý này cần được diễn giải đúng. Cơ sở hành vi không tự làm rỗng tập phương án
\emph{không phải vì nó dễ dãi}, mà vì nó đo so với một mốc mà nhà điều hành kiểm soát được. ``Không
làm gì'' theo định nghĩa là ``không gây ra gì''. Định lý \emph{không} nói hệ thống dùng cơ sở hành
vi không bao giờ dừng: tập phương án vẫn có thể rỗng vì một kiểm tra khác (mục \ref{muc:hocong}).

\begin{dinhly}[Cơ sở kết quả có thể vô nghiệm]
\label{dl:ketqua}
Dưới cơ sở kết quả, tồn tại trạng thái mà tập phương án của một xe bị làm rỗng.
\end{dinhly}

\begin{proof}
Với phương án không làm gì, $B = |p - v| = |e - v| = d$. Nếu sổ cái tích lũy cộng với $d$ vượt
$\beta$ thì ngay cả phương án không làm gì cũng bị loại. Khi mọi phương án khác cũng bị loại,
tập trở thành rỗng và hệ thống báo lỗi không còn phương án khả thi.
\end{proof}

\begin{dinhly}[Cơ sở kết quả không mở rộng tập quỹ đạo khả thi]
\label{dl:tapcon}
Giả sử điều kiện \emph{(H)}: tại mọi lần công bố, mốc ngoại sinh nằm gần phương án đang xét ít
nhất bằng mốc đã hứa, tức $|p - e| \le |p - v|$. Khi đó, với mọi trần $\beta$, mọi quỹ đạo thực
thi được dưới cơ sở kết quả đều thực thi được dưới cơ sở hành vi. Do đó số hành khách được phục
vụ bởi \emph{chính sách tối ưu} dưới cơ sở kết quả không lớn hơn số hành khách được phục vụ bởi
chính sách tối ưu dưới cơ sở hành vi.
\end{dinhly}

\begin{proof}
Điều kiện cần dùng là $A \le B$ tại mọi lần công bố. Vì $A = |p-e|$ và $B = |p-v|$, ta có

\begin{equation}
A \le B \iff |p - e| \le |p - v|,
\label{eq:dieukien}
\end{equation}

tức \emph{mốc ngoại sinh $e$ nằm gần phương án $p$ ít nhất bằng mốc đã hứa $v$}. Đây chính là
điều kiện (H) phát biểu ở dưới; giả thiết nó đúng trong phần chứng minh này.

Chứng minh bằng quy nạp theo số lần công bố. Ký hiệu $L^{A}_k$, $L^{B}_k$ là sổ cái tích lũy
của hai cơ sở sau $k$ lần công bố; ban đầu $L^{A}_0 = L^{B}_0 = 0$. Giả sử $L^{A}_k \le L^{B}_k$.
Nếu một phương án được cơ sở kết quả chấp nhận, tức $L^{B}_k + B \le \beta$, thì theo (H)
$$L^{A}_k + A \;\le\; L^{B}_k + B \;\le\; \beta,$$
nên phương án đó cũng được cơ sở hành vi chấp nhận, và bất đẳng thức giữa hai sổ cái được bảo
toàn sang bước $k+1$. Vậy mọi chuỗi quyết định hợp lệ dưới cơ sở kết quả cũng hợp lệ dưới cơ sở
hành vi, tức tập quỹ đạo khả thi của cơ sở kết quả là \emph{tập con} của tập quỹ đạo khả thi của
cơ sở hành vi. Lấy cận trên trên hai tập lồng nhau cho ngay kết luận.
\end{proof}

\paragraph{Điều kiện (H) và mức độ nó đúng.} Điều kiện \eqref{eq:dieukien} hỏng đúng khi trôi dạt
\emph{thuận chiều} đã đẩy mốc ngoại sinh ra xa phương án hơn so với mốc đã hứa --- nói cách khác,
khi trôi dạt thuận chiều tình cờ hấp thụ một phần cú đẩy của nhà điều hành, làm hành khách nhìn
thấy ít hơn lượng mà nhà điều hành thực sự gây ra. Đo trên các lần công bố thật của panel:

\begin{table}[htbp]
\centering
\caption{Mức độ đúng của điều kiện (H), tính trên từng lần công bố}
\label{bang:dieukien-h}
\small
\begin{tabular}{@{}lrr@{}}
\toprule
& \textbf{B1} & \textbf{C1}\\
\midrule
$A < B$ & 52.345 \;($96{,}38\%$) & 43.958 \;($96{,}52\%$)\\
$A = B$ & 1.738 \;($3{,}20\%$) & 1.581 \;($3{,}47\%$)\\
\midrule
\textbf{(H) thỏa} $A \le B$ & \textbf{54.083 \;($99{,}58\%$)} & \textbf{45.539 \;($99{,}99\%$)}\\
\textbf{(H) hỏng} $A > B$ & \textbf{228 \;($0{,}42\%$)} & \textbf{3 \;($0{,}007\%$)}\\
\quad vượt nhiều nhất & 690 ms & 104 ms\\
\bottomrule
\end{tabular}
\end{table}

Tập các lần công bố làm hỏng (H) \emph{trùng khít} với tập làm hỏng dấu bằng trong
\eqref{eq:phanra} --- cùng 228 bản ghi ở B1, cùng 3 bản ghi ở C1. Điều này là tất yếu: cả hai
đều là hệ quả của việc $e$ vượt qua $v$ theo chiều ngược với $p$.

Phần vi phạm cộng dồn bị chặn bởi tổng độ hụt ở bảng \ref{bang:khit-tung-lan}, tức 53.530 ms
trên toàn panel B1 --- $0{,}0748\%$ của $A+d$, và nhỏ hơn hai bậc độ lớn so với sàn nhiễu diễn
giải. Do đó Định lý \ref{dl:tapcon} được phát biểu như một kết quả \emph{có điều kiện}, với điều
kiện được đo chứ không được giả định. Cần nói rõ hai giới hạn của phép đo: (H) được đo trên các lần
công bố \emph{đã được chọn} của B1 và C1, không phải trên mọi phương án mà cơ sở kết quả có thể chấp
nhận; và vì nó hỏng 228 lần ở B1, giả thiết của định lý \emph{không} được thỏa nguyên vẹn trên
panel. Định lý là một mệnh đề về trường hợp (H) đúng, không phải một kết luận về panel này.

\paragraph{Phạm vi của Định lý \ref{dl:tapcon}, và điều nó \emph{không} nói.} Định lý phát biểu
về \emph{tập quỹ đạo khả thi} và về \emph{chính sách tối ưu}. Nó \emph{không} kéo theo rằng bộ
điều phối cận thị đang được cài đặt --- giải bài toán chọn phương án theo một hàm mục tiêu từ điển
ở từng kỷ nguyên, không nhìn trước các kỷ nguyên sau ---
sẽ phục vụ nhiều khách hơn dưới cơ sở hành vi trên \emph{một} quỹ đạo cụ thể. Lý do là mỗi lựa
chọn làm thay đổi trạng thái của các bước sau, nên một tập lựa chọn hẹp hơn tại mỗi trạng thái
vẫn có thể dẫn tới một quỹ đạo tình cờ tốt hơn. Bản thân thực nghiệm của khóa luận đã ghi nhận
hiện tượng này ở chiều khác: tỉ lệ phục vụ vi phạm tính đơn điệu ở 8 trên 20 ô khi siết trần
(mục \ref{muc:dovung}). Vì vậy phần thực nghiệm ở mục \ref{muc:haicoso} được trình bày như một
phép kiểm chứng độc lập, không phải như một hệ quả hiển nhiên của định lý.

Trong phạm vi đó, Định lý \ref{dl:tapcon} loại bỏ khả năng cơ sở kết quả \emph{về nguyên tắc}
vượt trội: khi (H) đúng, không có chính sách nào dưới cơ sở kết quả làm được điều mà một chính
sách dưới cơ sở hành vi không làm được. Kết quả trung tâm của chương lại nằm ở mục
\ref{muc:hocong}, nơi hai cơ sở này được đặt cạnh các ràng buộc lời hứa của văn liệu.

\begin{hequa}[Dải khác biệt tại một lần công bố]
\label{hq:dai}
Xét một lần công bố mà ngay trước đó hai sổ cái bằng nhau, $L^{A} = L^{B} = L$. Một phương án được
cơ sở hành vi giữ lại nhưng bị cơ sở kết quả loại khi và chỉ khi $L + A \le \beta < L + B$, tức
$L + A$ rơi vào dải $(\beta - (B - A),\, \beta]$. Vì $B \le A + d$, bề rộng dải không quá $d$ của
lần công bố đó, và đúng bằng $d$ khi dấu bằng trong \eqref{eq:phanra} xảy ra; khi (H) hỏng
($B < A$) dải rỗng.
\end{hequa}

Hệ quả này chỉ nói về \emph{một} lần công bố với hai sổ cái bằng nhau. Trên một quỹ đạo, hai sổ
cái tách nhau dần: sổ cái kết quả cộng thêm mọi trôi dạt đã xảy ra, nên độ chênh giữa hai cơ sở tăng
theo tổng trôi dạt tích lũy chứ không chỉ theo $d$ của một lần --- đúng điều phần (b) của Định lý
\ref{dl:hocong} dưới đây cho thấy, với ngưỡng an toàn của cơ sở kết quả bằng tổng trôi dạt. Nới trần
không làm dải của một lần công bố rộng ra, nhưng làm cho sổ cái kết quả khó chạm trần hơn; dưới (H)
và xét chính sách tối ưu, không tồn tại $\beta$ nào mà cơ sở kết quả phục vụ nhiều hơn. Trên panel
tĩnh, nới trần từ 30 lên 120 giây vẫn để 4 trên 20 ô vô nghiệm dưới cơ sở kết quả, chỉ là một tập ô
khác (bảng \ref{bang:doicho}).

\subsection{Khung so sánh bốn họ cổng}
\label{muc:hocong}

Ba định lý trên so hai cơ sở tính của \emph{cùng một} cơ chế. Các công trình gần nhất lại dùng
những ràng buộc lời hứa có hình thức khác: Schulz và Pfeiffer \cite{schulz2026} đặt một hạn chót
tính từ lời hứa đầu tiên; một cửa sổ quanh lời hứa đầu gần với ``hộp tĩnh'' gồm cửa sổ thời gian
và thời gian đi tối đa của Sharif Azadeh và cộng sự \cite{sharifazadeh2022}. Mục này đặt tất cả
vào cùng một khung để so được, và rút ra kết quả trung tâm của chương.

\paragraph{Đại lượng có dấu và đại lượng cộng dồn.} Xét một hành khách và tọa độ giờ trả. Tại mỗi
lần công bố $t$ đặt
\begin{equation}
x_t = e_t - v_t, \qquad c_t = p_t - e_t, \qquad z_t = p_t - v_t,
\end{equation}
khi đó $d_t = |x_t|$, $A_t = |c_t|$, $B_t = |z_t|$, và đồng nhất \emph{có dấu} $z_t = x_t + c_t$
luôn đúng --- bất đẳng thức \eqref{eq:phanra} chỉ xuất hiện khi lấy trị tuyệt đối. Gọi $p_0$ là
lời hứa \emph{đầu tiên} của hành khách. Cộng dồn trên cùng lưới công bố:
\begin{gather}
C = \sum_t A_t, \qquad X = \sum_t d_t, \qquad V = \sum_t B_t, \qquad E = |p_T - p_0|,\\
M = \max_t |p_t - p_0|, \qquad M^{+} = \max_t \max(0,\, p_t - p_0).
\end{gather}
$E$ là độ lệch hai đầu, $M^{+}$ là phần \emph{muộn} cực đại so với lời hứa đầu --- đúng đại lượng mà
một hạn chót kiểu Schulz--Pfeiffer chặn. Trong mục này $\kappa$ ký hiệu trần chung của một họ ràng
buộc; với hai sổ cái của khóa luận, $\kappa$ chính là trần $\beta$.

\begin{menhde}[Chuỗi đại lượng]
\label{md:chuoi}
Nếu chuỗi công bố liên tục, tức mỗi lần công bố so với lần ngay trước nó ($v_t = p_{t-1}$), thì
\[
\max(0,\, p_T - p_0) \;\le\; M^{+} \;\le\; M \;\le\; V \;\le\; C + X
\qquad\text{và}\qquad E \le M .
\]
\end{menhde}

\begin{proof}
Bất đẳng thức đầu và $E \le M$ đúng vì giá trị tại lần công bố cuối không vượt cực đại trên mọi
lần công bố; $M^{+} \le M$ vì $\max(0,\cdot) \le |\cdot|$. Với chuỗi liên tục,
$p_t - p_0 = \sum_{s \le t} z_s$, nên $|p_t - p_0| \le \sum_{s \le t} |z_s| \le V$. Cuối cùng
$|z_t| \le |c_t| + |x_t|$ ở từng lần công bố, cộng dồn lại được $V \le C + X$.
\end{proof}

Bản thân chuỗi bất đẳng thức này là kiến thức cơ bản, \emph{không} phải đóng góp; nó chỉ là công
cụ cho định lý dưới đây.

\paragraph{Bốn họ cổng.} Bảng \ref{bang:hocong} xếp bốn cách ràng buộc lời hứa theo đại lượng mà
chúng kiểm ở mỗi quyết định. (Schulz và Pfeiffer không công bố lại lời hứa --- khách chỉ được báo
giờ tới và hạn chót một lần --- nhưng hạn chót được kiểm ở mọi lần chèn khách mới về sau.)

\begin{table}[htbp]
\centering
\caption{Bốn họ cổng cam kết với trần $\kappa$}
\label{bang:hocong}
\small
\begin{tabular}{@{}llll@{}}
\toprule
\textbf{Họ} & \textbf{Kiểm ở mỗi quyết định} & \textbf{Tên gọi} & \textbf{Ví dụ}\\
\midrule
$C$ & $\sum_t A_t \le \kappa$ & theo hành vi & cơ sở hành vi của khóa luận\\
$M^{+}$ & $p_t - p_0 \le \kappa$ & hạn chót một chiều & Schulz và Pfeiffer \cite{schulz2026}\\
$M$ & $|p_t - p_0| \le \kappa$ & hạn chót hai chiều & cửa sổ quanh lời hứa đầu\\
$V$ & $\sum_t B_t \le \kappa$ & theo kết quả & cơ sở kết quả của khóa luận\\
\bottomrule
\end{tabular}
\end{table}

\paragraph{Phạm vi của định lý và giả thiết (N).} Định lý dưới đây chỉ nói về \emph{phần cam kết}
của bộ lọc. Tập phương án của một xe vẫn có thể rỗng, hoặc lần chạy vẫn có thể dừng, vì những ràng
buộc \emph{không} thuộc cơ chế cam kết: cấu trúc lộ trình, sức chứa, và nhất là bất biến vòng đời
khi xe đón muộn thật. Điều này đã xảy ra: dưới trôi dạt đo thật của cao điểm sáng (mục
\ref{muc:troidat}), chính cơ sở hành vi với trần 30 giây dừng ở 11 trên 20 ô, và cả 11 lần dừng
đều do bất biến cửa sổ đón của tầng miền (mục \ref{muc:batbien}) chứ không do cam kết. Cơ chế cụ
thể: với phương án không làm gì --- và \emph{chỉ} với nó --- bộ sinh nới giới hạn dịch vụ của các
khách trên tuyến đang chạy tới đúng mức mà chính tuyến đó đã vượt vì trôi dạt
(\texttt{ServiceQualityAllowance}; mọi phương án có thay đổi đều bị kiểm theo giới hạn gốc), nên
phương án này không bị loại vì cửa sổ; lần chạy dừng muộn hơn, khi xe thật sự tới điểm đón sau giờ đón muộn nhất
và thao tác lên xe bị tầng miền từ chối (\texttt{RideRequest.Board}). Gọi \textbf{(N)} là giả
thiết: \emph{tại kỷ nguyên đang xét, phương án không làm gì của xe đã qua mọi kiểm tra không thuộc
cam kết.}

\paragraph{Quỹ đạo dự phòng.} Xét quỹ đạo mà ở mọi kỷ nguyên xe giữ nguyên tuyến. Trên đó
$p_t = e_t$, nên $c_t = 0$ với mọi $t$, và
\[
C_{\text{noop}} = 0, \qquad z_t = x_t, \qquad V_{\text{noop}} = X_{\text{noop}} ,
\]
trong đó chỉ số ``noop'' chỉ đại lượng tính dọc quỹ đạo dự phòng (khác với quỹ đạo thật đã chạy).

\begin{dinhly}[Thứ tự khả thi của các họ cổng]
\label{dl:hocong}
Giả sử (N). Xét một hành khách.
\begin{enumerate}[label=(\alph*),leftmargin=1cm,itemsep=2pt,topsep=3pt]
\item \emph{Mọi quỹ đạo.} Với mọi $\kappa \ge 0$ và trên mọi quỹ đạo, cổng $C$ không bao giờ loại
phương án không làm gì. Do đó cổng hành vi \emph{không bao giờ là nguyên nhân} làm rỗng tập phương
án của một xe.
\item \emph{Quỹ đạo dự phòng.} Trên quỹ đạo dự phòng với chuỗi công bố liên tục, ngưỡng nhỏ nhất
bảo đảm phương án không làm gì không bị loại ở mọi kỷ nguyên là
\[
\kappa^{*}_{C} = 0 \;\le\; \kappa^{*}_{M^{+}} = M^{+}_{\text{noop}} \;\le\;
\kappa^{*}_{M} = M_{\text{noop}} \;\le\; \kappa^{*}_{V} = X_{\text{noop}} ,
\]
tức ba họ còn lại có thể loại phương án không làm gì, và ngưỡng an toàn của chúng tăng theo độ
trôi dạt --- với $M^{+}$ chỉ theo phần trôi \emph{bất lợi} (Hệ quả \ref{hq:chieutroi}).
\end{enumerate}
\end{dinhly}

\begin{proof}
(a) Sổ cái trước kỷ nguyên $t$ thỏa $C_{t-1} \le \kappa$ vì lần công bố trước đã qua cổng. Phương
án không làm gì có $c_t = 0$ nên $C_t = C_{t-1} \le \kappa$: cổng không loại nó, và theo (N) nó
cũng qua mọi kiểm tra khác, nên tập phương án không bị làm rỗng \emph{vì cổng}. (b) Phương án không
làm gì bị loại ngay tại lần công bố đầu tiên mà đại lượng được kiểm vượt $\kappa$, nên ngưỡng nhỏ
nhất đúng bằng giá trị lớn nhất của đại lượng đó dọc quỹ đạo dự phòng (với $V$ là giá trị cuối, vì
$V$ cộng dồn đơn điệu). Thứ tự giữa các ngưỡng là Mệnh đề \ref{md:chuoi} áp cho quỹ đạo có
$C = 0$; bước $M \le V$ cần chuỗi công bố liên tục.
\end{proof}

Phần (a) chính là Định lý \ref{dl:hanhvi} phát biểu lại trong khung chung; phần mới là việc đặt ba họ còn lại
--- trong đó có cơ chế của văn liệu --- vào cùng thước so. Chứng minh là \emph{sơ cấp}, và khóa
luận nói thẳng như vậy: giá trị của định lý không nằm ở độ khó mà ở hai chỗ --- đặt các họ cổng
trong văn liệu vào một khung để so được, và kiểm chứng rằng ranh giới nó vạch ra đúng với hệ thật.

Câu phát biểu chính xác là: \emph{cổng hành vi không bao giờ là nguyên nhân khiến tập phương án
rỗng}. Nói chặt hơn, phần (a) chứng minh rằng cổng $C$ không bao giờ loại phương án không làm gì;
``không bao giờ là nguyên nhân'' là cách diễn đạt điều đó: nếu tập phương án rỗng thì phương án
không làm gì đã bị một kiểm tra khác loại. Câu ``hệ thống dùng cổng hành vi không bao giờ kẹt'' là
\emph{sai} --- 11 trên 20 ô ở trên là phản ví dụ.

Tính ``duy nhất'' cũng phải đọc đúng phạm vi. Trong bốn họ, $C$ là họ \emph{duy nhất an toàn bất kể
chiều trôi}: hạn chót một chiều $M^{+}$ cũng không loại phương án không làm gì khi trôi hoàn toàn
thuận lợi (Hệ quả \ref{hq:chieutroi}), nhưng có thể loại nó khi trôi bất lợi.

\begin{hequa}[Khi thời gian đi tất định]
\label{hq:tatdinh}
Nếu thời điểm tới của tuyến giữ nguyên không trôi --- thời gian di chuyển và thời gian thực thi
đều tất định --- và chuỗi công bố liên tục, thì $x_t \equiv 0$, do đó
$M^{+}_{\text{noop}} = M_{\text{noop}} = X_{\text{noop}} = 0$ và cả bốn họ đều không bao giờ loại
phương án không làm gì, với mọi $\kappa \ge 0$.
\end{hequa}

Trong mô hình của Schulz và Pfeiffer, thời gian đi giữa hai điểm dừng là một giá trị cố định cho
mỗi cặp (mục 3.1 của bài), và bài khẳng định khách đã được nhận thì không bao giờ bị từ chối về
sau. Hệ quả \ref{hq:tatdinh} nhất quán với điều đó: trong một mô hình thời gian đi tất định, hạn
chót tính từ lời hứa đầu không thể tự nó làm rỗng tập phương án. Tính vô nghiệm chỉ xuất hiện khi
lời hứa của tuyến giữ nguyên bị trôi --- điều nằm ngoài mô hình của họ.

\begin{hequa}[Phụ thuộc chiều trôi]
\label{hq:chieutroi}
Trên quỹ đạo dự phòng với chuỗi công bố liên tục, $p_t - p_0 = \sum_{s \le t} x_s$, nên
\[
M^{+}_{\text{noop}} = \max_t \Big(\sum_{s \le t} x_s\Big)^{+}, \qquad
M_{\text{noop}} = \max_t \Big|\sum_{s \le t} x_s\Big|, \qquad
V_{\text{noop}} = \sum_t |x_t| .
\]
Nếu trôi \emph{bất lợi} (mọi $x_s \ge 0$, tắc đường tăng) thì
$M^{+}_{\text{noop}} = M_{\text{noop}} = V_{\text{noop}}$: ba họ $M^{+}$, $M$, $V$ cùng loại phương
án không làm gì ở cùng một ngưỡng. Nếu trôi \emph{thuận lợi} (mọi $x_s \le 0$, tắc đường giảm) thì
$M^{+}_{\text{noop}} = 0$: hạn chót một chiều không bao giờ loại phương án không làm gì, trong khi
$M_{\text{noop}} = V_{\text{noop}}$ vẫn có thể lớn, nên hạn chót hai chiều và cổng kết quả vẫn có
thể loại nó. Cổng hành vi không loại nó ở cả hai chiều.
\end{hequa}

Loại phương án không làm gì chưa phải là làm rỗng tập phương án. Với bộ sinh chèn thuần ($c \ge 0$)
và trôi bất lợi, mọi phương án khác đều muộn ít nhất bằng phương án không làm gì, nên chúng cũng bị
loại và tập phương án rỗng. Với trôi thuận lợi, một phương án làm muộn lại có thể lọt vào cửa sổ
của hạn chót hai chiều hay dưới trần kết quả; khi đó phương án không làm gì bị loại trong khi phương
án khác còn --- trạng thái mà mô hình chọn của hệ này, vốn đòi đúng một phương án không làm gì mỗi
xe, cũng không giải được.

Hệ quả này là chỗ làm cho định lý \emph{kiểm được} chứ không chỉ đúng. Một nhận định thô kiểu
``hạn chót hỏng khi có trôi dạt'' không phân biệt được hai chiều; Định lý \ref{dl:hocong} cùng Hệ
quả \ref{hq:chieutroi} đoán \emph{họ nào} hỏng theo \emph{chiều} trôi. Trên hệ thật $x_s$ không
thuần một dấu --- nó còn gồm tiến độ thực và độ trễ thực thi của xe --- nên các đẳng thức trên chỉ
đúng gần đúng, và mọi dự đoán phải tính từ quỹ đạo đo được của từng khách chứ không từ công thức.

Hệ quả cũng có một mặt trái phải được báo cáo. Dưới hạn chót một chiều và trôi thuận lợi, một
phương án chèn khách làm lời hứa muộn đi ($c_t > 0$) vẫn qua cổng khi $e_t - p_0$ đủ âm --- tức nhà
điều hành được \emph{tiêu phần trôi thuận lợi} để dời khách, trong khi cổng hành vi tính đủ $c_t$.
Vì vậy trong những khung giờ tắc đường giảm, hạn chót một chiều có thể phục vụ nhiều khách hơn trần
hành vi cùng mức, đổi lại khách thấy lời hứa đi sớm rồi muộn lại. Đây là một dự đoán phụ, chỉ nói
điều \emph{có thể} xảy ra; trên 8 ô khung tối của mục \ref{muc:tang2} nó \emph{không} được quan sát
(hạn chót 30 giây nhận thêm đúng một khách ở một ô, bằng nhau ở bảy ô). Nếu có, nó sẽ là một đánh đổi
giữa hai họ chứ không phải một ưu thế một chiều.

\subsection{Khi nào chặn theo chuỗi khác chặn hai đầu}
\label{muc:dondieu}

\begin{menhde}[Bộ sinh chèn thuần]
\label{md:dondieu}
Xét một chuỗi công bố liên tục. Nếu mọi phương án được chọn đều có $p \ge e$ với mọi hành khách
cũ, tức $c_t \ge 0$ ở mọi lần công bố, thì $C = \sum_t c_t$. Nếu thêm $x_t \equiv 0$ thì
$C = V = p_T - p_0 = M^{+} = M$: trần theo chuỗi, trần theo hai đầu và trần theo hành vi là một.
Khi $x_t \not\equiv 0$, kết luận thứ hai không còn đúng.
\end{menhde}

\begin{proof}
$c_t \ge 0$ cho $A_t = c_t$. Khi $x_t \equiv 0$ thì $z_t = c_t \ge 0$; với chuỗi liên tục
$p_t - p_0 = \sum_{s \le t} z_s$, nên $p_t$ đơn điệu không giảm, $p_T - p_0 = \sum_t z_t = V = C$,
và cực đại của $p_t - p_0$ đạt ở lần công bố cuối.
\end{proof}

Giả thiết thứ nhất của mệnh đề, $c_t \ge 0$, đúng trên hệ này ở mức đo được: dấu của $p - e$ dương
ở cả $479$ lần công bố có $A > 0$ của B1 (mục \ref{muc:dinhnghia-bamoc}), vì bộ sinh chỉ chèn thêm.
Do đó \emph{với phần do quyết định}, chặn theo chuỗi và chặn theo độ dời ròng là một: $C = \sum_t c_t$.
Giả thiết thứ hai, $x_t \equiv 0$, thì \emph{không} đúng: ngay trên panel tĩnh, tổng độ lớn phần
thay đổi khi giữ tuyến, $\sum_t |x_t|$, là $6.719$ giây ở B1 và $6.695$ giây ở C1, và nó đi cả hai
chiều gần như cân bằng (ở C1: $3.245$ giây muộn đi, $3.448$ giây sớm lên). Vì vậy phải tách hai
phát biểu.
\begin{itemize}[leftmargin=1.2cm,itemsep=2pt,topsep=3pt]
\item \emph{Với thành phần hành vi}, trên panel tĩnh, số khách vượt một ngưỡng tính theo độ trễ hai
đầu $D^{+} = \max(0, p_T - p_0)$ và tính theo chuỗi hành vi $C$ khác nhau không quá một khách trên
$1.735$ ở mọi ngưỡng 30, 60, 120, 300 giây. Chặn chuỗi hành vi hay chặn hai đầu gần như là một, và
khóa luận \emph{không} dùng điều đó làm điểm phân biệt.
\item \emph{Với chuỗi khách nhìn thấy} $V$, chuỗi và hai đầu tách nhau rõ: dưới C1, tổng $V$ là
$6.822$ giây trong khi tổng $D^{+}$ chỉ $126$ giây. $V$ cộng độ lớn từng bước còn $D^{+}$ chỉ nhìn
độ dời ròng dương, nên các dịch chuyển hai chiều vừa nêu triệt tiêu nhau ở hai đầu. Đây là nguồn
gốc của khoảng cách 54 lần ở mục
\ref{muc:doichieu}; và vì tổng độ lớn phần thay đổi khi giữ tuyến ở C1 ($6.695$ giây) gần bằng
$V$, phần bị thước hai đầu bỏ sót gần như toàn là thành phần mà cổng hành vi cố ý không tính.
\end{itemize}
Dưới trôi dạt thật, $D^{+}$ tách khỏi $C$ ngay cả trên cùng một quỹ đạo: ở $k = 1$, cực đại $D^{+}$
theo ô lên tới 102 giây trong khi tổng thành phần hành vi trên chín ô sống sót chỉ 27 giây. Chính
khoảng tách đó là thứ Định lý \ref{dl:hocong} khai thác.

\subsection{Các phản ví dụ}

Bảng \ref{bang:phanvidu} gom các tình huống nhỏ minh họa ranh giới của từng đại lượng. Đơn vị là
giây; CE-1 đến CE-5 là ví dụ dựng tay, CE-6 là một lần công bố thật.

\begin{table}[htbp]
\centering
\caption{Phản ví dụ cho các đại lượng của khung bốn họ cổng}
\label{bang:phanvidu}
\small
\begin{tabular}{@{}lp{3.9cm}p{3.1cm}p{5.2cm}@{}}
\toprule
& \textbf{Tình huống} & \textbf{Kết quả} & \textbf{Minh họa điều gì}\\
\midrule
CE-1 & $v = 600,\ e = 600,\ p = 660$ & $x = 0,\ c = 60,\ z = 60$ & Quyết định thuần gây dời.\\
CE-2 & $v = 600,\ e = 660,\ p = 660$ & $x = 60,\ c = 0,\ z = 60$ & Giữ tuyến mà lời hứa vẫn xấu
đi: trần hành vi \emph{không} chặn cái khách thấy.\\
CE-3 & $v = 600,\ e = 720,\ p = 600$ & $C = 120,\ V = 0$ & Sửa để giữ lời hứa bị trần hành vi
phạt; chưa từng xảy ra trên panel ($0/479$).\\
CE-4 & $p_0 = 600$; công bố $660$, rồi $600$ & $E = 0,\ V = 120$ & Thước hai đầu mù với dao
động.\\
CE-5 & $p_0 = 600$; công bố $660$, rồi $720$ & $E = 120,\ V = 120$ & Đơn điệu thì chuỗi bằng hai
đầu.\\
CE-6 & ca thật (xem dưới) & $A = 90{,}244$; $d = 0{,}690$; $B = 89{,}554$ & Trôi thuận chiều hấp
thụ một phần cú đẩy; hụt đúng $2d = 1{,}380$.\\
\bottomrule
\end{tabular}
\end{table}

Dòng CE-6 là một lần công bố thật ở ô 14-r1 của nhánh B1, lúc thời điểm mô phỏng 00:06:36:
$v = $ 00:15:09,943, $e = $ 00:15:09,253 và $p = $ 00:16:39,497 (giờ:phút:giây mô phỏng). Tuyến giữ
nguyên đã sớm hơn lời hứa $0{,}690$ giây, rồi quyết định đẩy giờ trả muộn thêm $90{,}244$ giây;
khách chỉ thấy $89{,}554$ giây. Đây cũng là một lần (H) hỏng ($A > B$), đúng loại ghi nhận ở bảng
\ref{bang:dieukien-h}.

Các kiểm chứng thực nghiệm của khung này nằm ở mục \ref{muc:troidat}, \ref{muc:haunghiem},
\ref{muc:hoito} --- giải thích hậu nghiệm trên dữ liệu có sẵn --- và mục \ref{muc:tang2}, một thí
nghiệm thăm dò có tiền đăng ký: dự đoán theo từng ô được niêm phong trước khi chạy các nhánh được
dự đoán và đúng ở
$98/102$ cặp trên panel chính, $31/32$ và $40/40$ ở khung sáng và khung tối của một tập giữ lại.

\section{Các cơ chế bảo vệ bổ sung}

Ngoài trần ngân sách, hệ thống cài đặt ba họ cơ chế bảo vệ khác, chủ yếu nhắm vào giờ đón. Việc
đưa cả ba vào là để so sánh được với các cơ chế đã có trong văn liệu.

\textbf{Khóa cứng} cấm hoàn toàn việc thay đổi một chiều nào đó sau khi lời hứa được xác nhận.
Cấu hình chính khóa cứng ba chiều định danh --- xe, điểm đón, điểm trả.

\textbf{Hàng rào đóng băng} (freeze horizon) cấm thay đổi khi thời điểm đón còn cách hiện tại
dưới một ngưỡng cho trước. Đây là phiên bản điều phối xe của cơ chế đóng băng lịch sản xuất
trong mục 2.3. Thực nghiệm dùng hai ngưỡng: 300 và 600 giây.

\textbf{Bảo đảm một chiều} (ratchet) cho phép giờ đón \emph{tốt lên} nhưng cấm \emph{xấu đi}.
Đây chính là cơ chế mà Alonso-Mora và cộng sự \cite{alonsomora2017} sử dụng, được mô tả trong
bài của họ là cấm tăng thời gian chờ đã đạt được so với phân công trước. Việc cài đặt nó cho
phép khóa luận so sánh trực tiếp ``cách của văn liệu'' với ``cách của chúng tôi'' trên cùng dữ
liệu, và kết quả ở mục \ref{muc:ratchet} khá bất ngờ.

\section{Kiến trúc cài đặt}

\subsection{Phân tầng}

Hệ thống được cài đặt bằng C\# trên nền .NET, tổ chức thành chín dự án theo kiến trúc phân tầng
với chiều phụ thuộc một chiều từ ngoài vào trong. Bảng \ref{bang:kientruc} tóm tắt.

\begin{table}[htbp]
\centering
\caption{Các thành phần chính của hệ thống}
\label{bang:kientruc}
\small
\begin{tabular}{@{}p{4.4cm}rp{7.4cm}@{}}
\toprule
\textbf{Thành phần} & \textbf{Dòng} & \textbf{Vai trò}\\
\midrule
\texttt{Domain} & 4.866 & Bất biến nghiệp vụ: chính sách cam kết, sổ cái, lời hứa, kiểm tra ràng buộc vật lý.\\
\texttt{Application} & 4.023 & Vòng quyết định, bộ tính độ dời ba mốc, bộ kiểm tra cam kết.\\
\texttt{Algorithms} & 7.889 & Sinh phương án theo lối chèn, chứng chỉ độ trễ an toàn, các hồ sơ hàm mục tiêu.\\
\texttt{Solvers.OrTools} & 687 & Cầu nối tới bộ giải quy hoạch nguyên của OR-Tools.\\
\texttt{Contracts} & 9.332 & Giao thức trao đổi với bộ mô phỏng, mã hóa chuẩn tắc.\\
\texttt{Runner} & 6.295 & Tiến trình chạy, chuẩn tắc hóa trạng thái, ghi biên bản.\\
\texttt{Benchmarking} & 17.150 & Chuẩn hóa dữ liệu, dựng bộ dữ liệu thử, sinh chứng cứ.\\
\bottomrule
\end{tabular}
\end{table}

Bộ mô phỏng sử dụng là FleetPy phiên bản 1.0.2 \cite{fleetpy2022}, một nền mô phỏng mở cho các
dịch vụ vận tải theo yêu cầu. Hệ thống điều phối của khóa luận đóng vai trò bộ điều khiển đội
xe ngoài, trao đổi với FleetPy qua một giao thức nhị phân có kiểm tra băm ở từng khung tin.

\subsection{Sinh phương án theo lối chèn}

Bộ sinh phương án hoạt động theo nguyên tắc chèn: với mỗi xe và mỗi yêu cầu mới, nó thử chèn
cặp điểm đón--trả vào mọi vị trí hợp lệ của lộ trình hiện tại, giữ nguyên thứ tự các điểm đã có.
Đặc điểm này có hai hệ quả quan trọng đối với phần lý thuyết. Thứ nhất, nó là lý do
$\mathtt{decisionMovedEarlier} = 0$: chèn thêm chỉ đẩy lùi được giờ của khách cũ. Thứ hai, nó
giới hạn không gian tìm kiếm một cách có kiểm soát, cho phép khóa luận chứng minh rằng kết quả
không phải do tìm kiếm bị cắt ngắn --- điểm sẽ được kiểm chứng ở mục \ref{muc:bosinh}.

Hệ thống cũng cài đặt một cơ chế \emph{chứng chỉ độ trễ}: một phép tính bảo thủ chứng minh rằng
một độ trễ thuần túy là an toàn, dùng để dời phần chờ đã có sẵn tới một điểm chờ thực thi được
mà không phát sinh thêm độ trễ mới.

\subsection{Chứng cứ niêm phong}

Vì toàn bộ kết luận của khóa luận dựa trên so sánh giữa các cấu hình, tính toàn vẹn của dữ liệu
thực nghiệm là điều kiện tiên quyết. Hệ thống ghi ba loại chứng cứ cho mỗi lần chạy.

Thứ nhất là \emph{biên bản đầy đủ}: mọi khung tin trao đổi giữa bộ mô phỏng và bộ điều phối
được ghi ra dưới dạng NDJSON, mã hóa base64, kèm băm SHA-256 từng khung. Từ biên bản này có thể
dựng lại toàn bộ diễn biến, kể cả những đại lượng không nằm trong bản tóm tắt.

Thứ hai là \emph{băm ngữ nghĩa} (\texttt{semanticHash}) và \emph{tóm lược công bố}
(\texttt{publicationDigest}): hai giá trị băm tổng hợp trạng thái cuối và toàn bộ chuỗi lời hứa
đã công bố. Hai lần chạy khác nhau về hành vi chắc chắn khác nhau ở hai giá trị này. Cơ chế
này đã tỏ ra thiết yếu trong mục \ref{muc:ratchet}, nơi sáu cấu hình khác nhau cho ra đúng hai bộ kết quả ---
nếu không có băm phân biệt thì không thể loại trừ khả năng đó là lỗi sao chép cấu hình.

Thứ ba là \emph{băm nhị phân bộ chạy} (\texttt{runnerSha256}) và băm kho mã nguồn, ghim kết quả
vào đúng phiên bản đã sinh ra nó. Cơ chế này cũng chính là thứ đã chặn hai cấu hình trong mục
4.6 --- bộ chạy đã niêm phong ra đời trước khi tính năng ratchet vào mã nguồn nên từ chối nạp
cấu hình đó --- và buộc khóa luận phải tách thành một thí nghiệm phụ riêng thay vì trộn lẫn kết
quả từ hai phiên bản khác nhau.
